\section{Policy Gradient 的基本推导}

\subsection{从期望奖励到 log-derivative trick}

为了记号简洁，把整段 response 当作一个整体动作 $a$。那么
\[
\mathbb{E}[R] = \int p(s)\,\pi(a \mid s)\,R(s,a)\,ds\,da.
\]

对 policy 参数求梯度：
\[
\nabla \mathbb{E}[R]
= \int p(s)\,\nabla \pi(a \mid s)\,R(s,a)\,ds\,da
= \int p(s)\,\pi(a \mid s)\,\nabla \log \pi(a \mid s)\,R(s,a)\,ds\,da.
\]

于是得到熟悉的形式：
\[
\nabla \mathbb{E}[R] = \mathbb{E}\big[\nabla \log \pi(a \mid s)\,R(s,a)\big].
\]

如果把 response 展开成 token 序列 $a_{1:T}$，那 policy 其实在每一步都给出一个条件分布：
\[
\pi(a_{1:T}\mid s)=\prod_{t=1}^T \pi(a_t \mid s, a_{<t}),
\]
所以
\[
\log \pi(a_{1:T}\mid s)=\sum_{t=1}^T \log \pi(a_t \mid s, a_{<t}).
\]
这件事非常重要，因为它说明 \textbf{sequence-level 的 reward 最后会分摊到每一个 token 的 log prob 上}。虽然最后只给一个分数，但梯度会回流到整条生成轨迹。

\begin{importantbox}{policy gradient 的直觉}
\begin{itemize}
    \item 采到好样本，就提高它的概率；
    \item 采到差样本，就降低它的概率；
    \item 变化的幅度由 reward 决定。
\end{itemize}
\end{importantbox}


\begin{knowledgebox}{为什么要对 log prob 求梯度}
`log` 把乘法关系变成加法，梯度也更稳定。更重要的是，它让“提高一个 sample 的概率”变成一个直接可微的信号，便于把 reward 接到反向传播里。
\end{knowledgebox}


\subsection{naive policy gradient}

最朴素的做法就是：
\begin{enumerate}
    \item 采样一个 prompt $s$；
    \item 从当前 policy 采样 response $a$；
    \item 用 $\nabla \log \pi(a \mid s)\,R(s,a)$ 更新参数。
\end{enumerate}

这和 supervised fine-tuning 很像，只不过监督数据不是人工标注，而是模型自己采样出来后再按 reward 过滤和加权。

\begin{warningbox}{稀疏奖励会把更新打空}
如果 $R \in \{0,1\}$，多数样本 reward 都是 0，那么很多 batch 的梯度就是近似 0。模型如果一开始就很差，可能连“碰到正奖励样本”的机会都很少。
\end{warningbox}


